% META: {"id":"1SPE-TRIGO-ME-002","chapitre":"1SPE-TRIGONOMETRIE","type_objet":"methode","capacites_codes":["C2"],"methodes":["M2"],"sources_inspiration":[],"status":"approved","origin":"generated","status_history":["generated","audited","accepted","approved"],"release_acceptance":"RELEASE_OWNER_BATCH_ACCEPTANCE","acceptance_closure_digest":"sha256:9b3ccf9a81c5520fbb7e03b7b2d3e7bf2057a02834b6d908bf3be5f49f3b553f","acceptance_receipt_digest":"sha256:4fad86f0282e0fa4e0419cca1ad6379ae7af5a280c04a4a90880f90a1c044242"}
\begin{fichemethode}{M2}{Calculer un cosinus ou un sinus avec les valeurs remarquables et les symetries}

\quandUtiliser{%
  Utiliser cette méthode lorsqu'on doit calculer la valeur exacte de $\cos\alpha$ ou $\sin\alpha$
  pour un angle $\alpha$ qui n'est pas directement dans le tableau des valeurs remarquables
  mais qui s'y ramene par symetrie.
}

\pasApas{%
  \begin{enumerate}
    \item Déterminer la mesure principale de $\alpha$.
    \item Identifier a quel angle remarquable $\beta \in \{0, \frac{\pi}{6}, \frac{\pi}{4}, \frac{\pi}{3}, \frac{\pi}{2}\}$ on peut se ramener.
    \item Appliquer la formule de symetrie adaptee :
      \begin{itemize}
        \item Si $\alpha = -\beta$ : parite.
        \item Si $\alpha = \pi - \beta$ : supplementaire.
        \item Si $\alpha = \pi + \beta$ : anti-supplementaire.
        \item Si $\alpha = \frac{\pi}{2} - \beta$ : complementaire.
      \end{itemize}
    \item Remplacer par la valeur remarquable connue.
  \end{enumerate}
}

\exempleRedige{%
  Calculer $\cos\!\left(\dfrac{5\pi}{4}\right)$ et $\sin\!\left(\dfrac{5\pi}{4}\right)$.

  \medskip
  \commentaireMarge{Écrire $\frac{5\pi}{4} = \pi + \frac{\pi}{4}$.}%
  On a $\dfrac{5\pi}{4} = \pi + \dfrac{\pi}{4}$.

  \commentaireMarge{Appliquer $\cos(\pi + x) = -\cos x$ et $\sin(\pi + x) = -\sin x$.}%
  Donc :
  \[
    \cos\!\left(\frac{5\pi}{4}\right) = -\cos\!\left(\frac{\pi}{4}\right) = -\frac{\sqrt{2}}{2}, \qquad
    \sin\!\left(\frac{5\pi}{4}\right) = -\sin\!\left(\frac{\pi}{4}\right) = -\frac{\sqrt{2}}{2}.
  \]
}

\pieges{%
  \begin{itemize}
    \item \textbf{Piège 1.} Inverser $\cos\!\left(\dfrac{\pi}{3}\right) = \dfrac{1}{2}$ et $\cos\!\left(\dfrac{\pi}{6}\right) = \dfrac{\sqrt{3}}{2}$.
      Moyen mnemotechnique : quand l'angle augmente de $0$ a $\frac{\pi}{2}$, le cosinus diminue.
    \item \textbf{Piège 2.} Erreur de signe dans la formule du supplementaire :
      $\cos(\pi - x) = -\cos x$ (signe $-$) mais $\sin(\pi - x) = +\sin x$ (signe $+$).
  \end{itemize}
}

\verifier{%
  Vérifier que $\cos^2\alpha + \sin^2\alpha = 1$. Si ce n'est pas le cas, il y a une erreur de calcul.
}

\sEntrainer{\refExos{M2}}

\end{fichemethode}
